\subchapter{Restricting userland applications}{Objective: Learn
	how to restrict userspace attack surface using various tools: seccomp, SELinux,
	and systemd configuration}

To demonstrate possibilities to filter userspace operations, we will show how
various mechanisms can be used. First, we will use seccomp to restrict
accessible syscalls in a simple application.
Next, we will quickly manipulate SELinux contexts.
Finally, we will show how userland security features can be used directly from
Systemd.

\section{Preventing a simple application from executing shellcode using SECCOMP}

The \code{$HOME/__SESSION_NAME__-labs/userland/seccomp} folder contains the
code of a sample application.
This application does two things: it prints the content of some existing file,
then launches a shell.
Of course, launching the shell is an expected behavior here, but we are
actually trying to mimic a faulty application, where an attacker exploits a
breach to launch a shell.
To avoid that, we will demonstrate how \emph{seccomp} can be used to restrict
accessible syscalls.

First, you can have a look at the source code and compile the application.
Then, launch it and observe the current behaviour.

\begin{bashinput}
$ cd $HOME/__SESSION_NAME__-labs/userland/seccomp
$ . $HOME/__SESSION_NAME__-labs/sdk/environment-setup-cortexa55-oe-linux
$ make
$ ssh root@${BOARD} mkdir -p /usr/local/bin
$ scp seccomp-test root@${BOARD}:/usr/local/bin
$ ssh root@${BOARD} /usr/local/bin/seccomp-test
Showing content of "/proc/cmdline":
console=ttyLP0,115200 root=/dev/mmcblk1p2 rootwait rw selinux=1 enforcing=0
Running "/bin/bash":
\end{bashinput}

As you can see, after printing file content, a shell was started.
We really want to avoid that.

One solution is to use basic seccomp mechanism.
Remember, this will only allow processes to use very basic syscalls, such as
\texttt{read()} or \texttt{write()}.
Modify the \texttt{seccomp-test.c} source file, and uncomment the call to
\texttt{enable\_seccomp()}.
Rebuild, and relaunch.

\begin{bashinput}
$ make
$ scp seccomp-test root@${BOARD}:/usr/local/bin
\end{bashinput}

We do not get much output with SSH, so you should run it directly on the target:

\begin{bashinput}
$ ssh root@${BOARD} /usr/local/bin/seccomp-test
Killed
\end{bashinput}

As expected, the process gets killed.
However, what might not be expected is the content of \texttt{/proc/cmdline} is
no longer shown.
You can use \texttt{strace} to analyze what syscalls were used, and understand
this new behaviour.
On the board, you can run:

\begin{bashinput}
$ strace seccomp-test
...
seccomp(SECCOMP_SET_MODE_STRICT, 0, NULL) = 0
fstat(1,  <unfinished ...>)             = ?
+++ killed by SIGKILL +++
[1]    1202806 killed     strace seccomp-test
\end{bashinput}

As you can see, \texttt{printf()} did call \texttt{fstat()}.
It tends to do that only on the very first call, so this could be avoided by
calling \texttt{printf()} before enabling seccomp.
But we are later calling \texttt{close()}, which is also denied by seccomp.

As an alternative, we will use libseccomp to add our own custom filters.

Modify again the source code, defining the \texttt{USE\_SECCOMP\_FILTERS} macro:
the \texttt{enable\_seccomp()} will now use libseccomp to define our custom
filters.
Some filters were already added by the example code, but you will have to
complete this list.
Again, you can use \texttt{strace} to identify the correct syscalls.

Remember, our goal is to have an application that can successfully show the file
content, but cannot execute the shell.
After having added all needed filters, you should have the following behaviour:

\begin{bashinput}
$ seccomp-test
Showing content of "/proc/cmdline":
console=ttyLP0,115200 root=/dev/mmcblk1p2 rootwait rw selinux=1 enforcing=0
Running "/bin/bash":
Bad system call (core dumped)
\end{bashinput}

\section{Manipulating SELinux contexts}

To demonstrate SELinux use case, we will use a small TFTP server present on the
board.
This TFTP server has been misconfigured: instead of serving the content of
\texttt{/srv/tftp/}, it does serve the content of the whole root filesystem
(\texttt{/}).
We can create two files on the board, to demonstrate this behaviour.

\begin{bashinput}
# echo "Public data" > /srv/tftp/public
# echo "Very secret data" > /root/secret
# chmod 777 /root
\end{bashinput}

On the host computer, use a tftp client to download these files.
As the tftp server will connect to a random UDP port on the host, you might need
to disable some firewalling rules.

\begin{bashinput}
$ tftp $BOARD -vc get /srv/tftp/public
Connected to 192.168.0.168 (192.168.0.168), port 69
getting from 192.168.0.168:/srv/tftp/public to public [netascii]
Received 13 bytes in 0.0 seconds [3200 bit/s]
$ cat public
Public data
$ tftp $BOARD -vc get /root/secret
Connected to 192.168.0.168 (192.168.0.168), port 69
getting from 192.168.0.168:/root/secret to secret [netascii]
Received 18 bytes in 0.0 seconds [4270 bit/s]
$ cat secret
Very secred data
\end{bashinput}

If downloads are failing, you might need to have a look at the tftp server logs:

\begin{bashinput}
# journalctl -u atftpd -f
\end{bashinput}

Once we have this working setup, we can start experimenting with SELinux.
SELinux enforcement is disabled by default, thanks to the kernel command line.
This is a specific setup, made so we can run other labs without having SELinux
blocking anything.
The first step is then to enable SELinux until the next reboot:

\begin{bashinput}
# setenforce 1
\end{bashinput}

SELinux status can be shown with the \texttt{sestatus} command:

\begin{bashinput}
# sestatus
SELinux status:                 enabled
SELinuxfs mount:                /sys/fs/selinux
SELinux root directory:         /etc/selinux
Loaded policy name:             standard
Current mode:                   enforcing
Mode from config file:          enforcing
Policy MLS status:              disabled
Policy deny_unknown status:     allowed
Memory protection checking:     actual (secure)
Max kernel policy version:      34
\end{bashinput}

Let's manipulate a bit SELinux: first we will look for SELinux users.
You can both the list of users and the mapping between Linux and SELinux users.

\begin{bashinput}
# seinfo -u

Users: 7
   root
   staff_u
   sysadm_u
   system_u
   unconfined_u
   user_u
   xdm
# semanage login -l

Login Name                SELinux User

__default__               user_u
root                      root
\end{bashinput}

The first commands, shows existing SELinux users, the second one shows mapping
from Linux users.
We can see the \emph{root} user is mapped to the \emph{root} SELinux user, while
all others are mapped to \emph{user\_u}.

Similarly, we can show the corresponding roles:

\begin{bashinput}
# seinfo -u -x

Users: 7
   user root roles { staff_r sysadm_r system_r };
   user staff_u roles { staff_r sysadm_r };
   user sysadm_u roles sysadm_r;
   user system_u roles system_r;
   user unconfined_u roles { system_r unconfined_r };
   user user_u roles user_r;
   user xdm roles xdm_r;
\end{bashinput}

As our Linux root user is mapped to the SELinux root user, the corresponding
line is the first one. Keep in mind this mapping might be different on
production systems.

We can list which types can be accessed by \emph{staff\_r}, \emph{sysadm\_r} or
\emph{system\_r} roles.

\begin{bashinput}
# seinfo -rstaff_r -x

Roles: 1
   role staff_r types { auditadm_screen_t bluetooth_helper_t cdrecord_t chfn_t chkpwd_t chromium_naclhelper_t chromium_renderer_t chromium_sandbox_t chromium_t consoletype_t container_engine_t container_t crio_t cronjob_t crontab_t ddclient_t dhcpc_t dirmngr_t 
   ...

# seinfo -rsysadm_r -x

Roles: 1
   role sysadm_r types { acngtool_t acpi_t admin_crontab_t aide_t amanda_recover_t apt_t auditadm_screen_t auditctl_t backup_t bacula_admin_t bluetooth_helper_t bootloader_t calamaris_t cdrecord_t certbot_t certwatch_t cgclear_t checkpc_t checkpolicy_t chfn_t
   ...

# seinfo -rsystem_r -x

Roles: 1
   role system_r types { NetworkManager_t abrt_dump_oops_t abrt_handle_event_t abrt_helper_t abrt_retrace_coredump_t abrt_retrace_worker_t abrt_t abrt_upload_watch_t abrt_watch_log_t accountsd_t acct_t acngtool_t acpi_t acpid_t afs_bosserver_t afs_fsserver_t
   ...
\end{bashinput}

We can also observe the context of running processes, with \texttt{ps -Z}:

\begin{bashinput}
# ps -Z
  PID CONTEXT                          STAT COMMAND
    1 system_u:system_r:init_t         S    {systemd} /sbin/init
    2 system_u:system_r:kernel_t       SW   [kthreadd]
    3 system_u:system_r:kernel_t       SW   [pool_workqueue_]
    4 system_u:system_r:kernel_t       IW<  [kworker/R-rcu_g]
    5 system_u:system_r:kernel_t       IW<  [kworker/R-sync_]
    6 system_u:system_r:kernel_t       IW<  [kworker/R-kvfre]
    ...
  310 root:sysadm_r:sysadm_t           S    -sh
  338 system_u:system_r:kernel_t       IW<  [kworker/1:2H]
  371 system_u:system_r:kernel_t       IW   [kworker/u8:1-ev]
  386 system_u:system_r:kernel_t       IW   [kworker/u8:2-ev]
  389 root:sysadm_r:sysadm_t           R    {ps} /usr/bin/busybox.nosuid /usr/bin/ps -Z
\end{bashinput}

We now want to compare the SELinux context of the two files we created earlier.
This can be shown with \texttt{ls -Z}.

\begin{bashinput}
# ls -lZ /srv/tftp/public /root/secret
-rw-r--r--. 1 root root root:object_r:user_home_t 17 Mar 17 09:27 /root/secret
-rw-r--r--. 1 root root root:object_r:var_t       12 Mar 17 09:31 /srv/tftp/public
\end{bashinput}

As you can see, the two files were automatically labelled differently, based on
the folder they were created in.
As SELinux is now enabled, we can again try to dowload them using tftp.
You can run the same commands as previously on the host

\begin{bashinput}
$ rm -f public secret
$ tftp $BOARD -vc get /srv/tftp/public
Connected to 192.168.0.168 (192.168.0.168), port 69
getting from 192.168.0.168:/srv/tftp/public to public [netascii]
Received 13 bytes in 0.0 seconds [3200 bit/s]
$ cat public
Public data
$ tftp $BOARD -vc get /root/secret
Connected to 192.168.0.168 (192.168.0.168), port 69
getting from 192.168.0.168:/root/secret to secret [netascii]
Error code 1: File not found
\end{bashinput}

Thanks to the SELinux policies, the tftp server no longer has access to the
content of \texttt{/root}.

Let's pretend this access is legitimate.
We could either modify SELinux rules or the file context.
For demonstration purpose, we can try this later solution, changing the context
of \texttt{/root/secret} to be the same as \texttt{/srv/tftp/public}.

\begin{bashinput}
# chcon root:object_r:var_t /root/secret
# ls -lZ /root/secret
-rw-r--r--. 1 root root root:object_r:var_t 17 Mar 17 09:27 /root/secret
\end{bashinput}

And then try again to download the file:

\begin{bashinput}
$ tftp $BOARD -vc get /root/secret
Connected to 192.168.0.168 (192.168.0.168), port 69
getting from 192.168.0.168:/root/secret to secret [netascii]
Received 18 bytes in 0.0 seconds [4580 bit/s]
\end{bashinput}

The \texttt{restorecon} command can then be used to revert to a sane context:

\begin{bashinput}
# restorecon /root/secret
root@freiheit93:~# ls -lZ /root/secret
-rw-r--r--. 1 root root root:object_r:user_home_t 17 Mar 17 09:27 /root/secret
\end{bashinput}

\section{Using systemd to restrict a daemon's access to resources}

We now want to run an application with systemd, using systemd unit configuration
to filter the operations the application can make.

This application is again an artificial one, trying to mimic one doing some legit
operations, but also doing some unwanted ones, because of bugs or security
breaches exploited by an attacker.

As legitimate operation, the application does:
\begin{itemize}
  \item Adds some entropy to \texttt{/dev/urandom}: this does require the
    \texttt{CAP\_SYS\_ADMIN} capability.
    Note, this is mostly used for demonstration.
    The added entropy is based on current system time: this is not a
    particularly good entropy source and should not be used in real systems.
  \item Runs \texttt{ping -N}.
    This is again an option requiring the \texttt{CAP\_NET\_RAW} capability.
\end{itemize}

The application also simulates two bogus operations:
\begin{itemize}
	\item It writes some data to a file in \texttt{/var/cache}.
	\item It modifies the machine hostname.
\end{itemize}

The \code{$HOME/__SESSION_NAME__-labs/userland/systemd} contains both the
application source code and a sample systemd unit.
Cross-compile the application, and deploy it on the target, to the
\texttt{/usr/local/bin} folder; then deploy the service file to
\texttt{/etc/systemd/system/}.

\begin{bashinput}
$ cd $HOME/__SESSION_NAME__-labs/userland/systemd
$ make
$ scp systemd-test root@${BOARD}:/usr/local/bin
$ scp training-systemd-test.service root@${BOARD}:/etc/systemd/system/
\end{bashinput}

Now, start the service and monitor the log output.

\begin{bashinput}
# systemctl daemon-reload
# systemctl restart training-systemd-test --no-block
# journalctl -u training-systemd-test -f -n 50
Starting training-systemd-test.service - Security training systemd demo...
Adding entropy count
Writing data to "/var/cache/security_training_systemd"
Setting hostname to "changed-hostname"
hostname: changed-hostname
Executing ping command
PING google.com (172.217.22.206) 56(84) bytes of data.
64 bytes from muc11s01-in-f14.1e100.net (172.217.22.206): icmp_seq=1 ttl=111 time=14.1 ms
--- google.com ping statistics ---
1 packets transmitted, 1 received, 0% packet loss, time 0ms
rtt min/avg/max/mdev = 14.112/14.112/14.112/0.000 ms
training-systemd-test.service: Deactivated successfully.
Finished training-systemd-test.service - Security training systemd demo.
\end{bashinput}

Looking at the log, we have some good and bad news: adding entropy and running
ping did succeed.
But so did changing the hostname and writing data to
\texttt{/var/cache/security\_training\_systemd}.
We can verify this:

\begin{bashinput}
# cat /var/cache/security_training_systemd
Hello!
# cat /proc/sys/kernel/hostname
changed-hostname
\end{bashinput}

Let's revert these changes:

\begin{bashinput}
# rm /var/cache/security_training_systemd
# echo "myboard" > /proc/sys/kernel/hostname
\end{bashinput}

To prevent these changes, we can limit the permissions of the process launched
by systemd.
First, instead of running as root, we should be using an unprivileged user
account, thanks to the \texttt{User=} directive.
We can ask systemd to create an ephemeral one, with \texttt{DynamicUser=}.

Modify the unit file with \texttt{systemctl edit --full training-systemd-test}
to add following lines into the \texttt{[Service]} section.
Then restart the service.

\begin{verbatim}
User=someuser
DynamicUser=true
\end{verbatim}

\begin{bashinput}
# systemctl restart training-systemd-test --no-block
# journalctl -u training-systemd-test -f -n 50
Starting training-systemd-test.service - Security training systemd demo...
Writing data to "/var/cache/security_training_systemd"
Failed to open "/var/cache/security_training_systemd": Read-only file system
training-systemd-test.service: Main process exited, code=exited, status=1/FAILURE
training-systemd-test.service: Failed with result 'exit-code'.
Failed to start training-systemd-test.service - Security training systemd demo.
\end{bashinput}

The application can no longer write to \texttt{/var/cache/}, that's great, but
as a consequence it immediately stops.
We can ask systemd to mount a tmpfs on /var/cache/, so the application can still
write in there, without impacting the system.
Modify the service file, and add the following line:

\begin{verbatim}
TemporaryFileSystem=/var/cache:mode=0777
\end{verbatim}

Then restart the application, and look for logs.

\begin{bashinput}
# systemctl restart training-systemd-test --no-block
# journalctl -u training-systemd-test -f -n 50
Starting training-systemd-test.service - Security training systemd demo...
Writing data to "/var/cache/security_training_systemd"
Adding entropy count
Failed to add entropy: Operation not permitted
Setting hostname to "changed-hostname"
Failed to set hostname: Operation not permitted
hostname: myboard
Executing ping command
ping: socktype: SOCK_RAW
ping: socket: Operation not permitted
ping: => missing cap_net_raw+p capability or setuid?
training-systemd-test.service: Main process exited, code=exited, status=2/INVALIDARGUMENT
training-systemd-test.service: Failed with result 'exit-code'.
Failed to start training-systemd-test.service - Security training systemd demo.
\end{bashinput}

So now, the application can happily write to
\texttt{/var/cache/security\_training\_systemd}.
You can verify that this file was not really created on your system, but on some
tmpfs.
We can also see the hostname was not changed.
However, we now have two issues: adding entropy failed, and ping
complains: \texttt{WARNING: failed to set mark: 42: Operation not permitted}.
Both issues come from missing capabilities, as we are no longer running as
root.
We can modify the configuration to add the required ambient capabilities and
restart the service.

\begin{verbatim}
AmbientCapabilities=CAP_NET_RAW CAP_SYS_ADMIN
\end{verbatim}

\begin{bashinput}
# systemctl restart training-systemd-test --no-block
# journalctl -u training-systemd-test -f -n 50
Starting training-systemd-test.service - Security training systemd demo...
Writing data to "/var/cache/security_training_systemd"
Adding entropy count
Setting hostname to "changed-hostname"
hostname: changed-hostname
Executing ping command
PING google.com (172.217.22.78) 56(84) bytes of data.
64 bytes from tzpara-am-in-f14.1e100.net (172.217.22.78): icmp_seq=1 ttl=112 time=14.1 ms
--- google.com ping statistics ---
1 packets transmitted, 1 received, 0% packet loss, time 0ms
rtt min/avg/max/mdev = 14.122/14.122/14.122/0.000 ms
training-systemd-test.service: Deactivated successfully.
Finished training-systemd-test.service - Security training systemd demo.
\end{bashinput}

With the added capabilities, both the ping and entropy adding operation succeed,
but we also granted the possibility to change the hostname.
One solution is to use the \texttt{ProtectHostname=} option to forbid this.

\begin{verbatim}
ProtectHostname=on
\end{verbatim}

\begin{bashinput}
# systemctl restart training-systemd-test --no-block
# journalctl -u training-systemd-test -f -n 50
Starting training-systemd-test.service - Security training systemd demo...
Writing data to "/var/cache/security_training_systemd"
Adding entropy count
Setting hostname to "changed-hostname"
Failed to set hostname: Operation not permitted
hostname: myboard
Executing ping command
PING google.com (172.217.22.78) 56(84) bytes of data.
64 bytes from tzpara-am-in-f14.1e100.net (172.217.22.78): icmp_seq=1 ttl=112 time=13.8 ms
--- google.com ping statistics ---
1 packets transmitted, 1 received, 0% packet loss, time 0ms
rtt min/avg/max/mdev = 13.830/13.830/13.830/0.000 ms
training-systemd-test.service: Deactivated successfully.
Finished training-systemd-test.service - Security training systemd demo.
\end{bashinput}

Alternatively, one could use seccomp to prevent call to \texttt{sethostname()}.
Systemd \texttt{SystemCallFilter=} configuration allows to do so easily.

\begin{verbatim}
SystemCallFilter=~sethostname
\end{verbatim}

An interesting difference is the process is killed, except if we also use
\texttt{SystemCallErrorNumber=}.

\begin{bashinput}
# systemctl restart training-systemd-test --no-block
# journalctl -u training-systemd-test -f -n 50
Starting training-systemd-test.service - Security training systemd demo...
Writing data to "/var/cache/security_training_systemd"
Adding entropy count
Setting hostname to "changed-hostname"
training-systemd-test.service: Main process exited, code=killed, status=31/SYS
training-systemd-test.service: Failed with result 'signal'.
Failed to start training-systemd-test.service - Security training systemd demo.
\end{bashinput}
